% =====================================================================
% CAPÍTULO 7: PLANIFICACIÓN, CRONOGRAMA Y CONCLUSIONES
% =====================================================================
\chapter{Planificación, Cronograma y Conclusiones}
\label{ch:cronograma}

\section{Cronograma General del Proyecto}

El proyecto se planifica en 24 semanas, distribuidas en ocho fases alineadas con
el Capítulo~\ref{ch:metodologia}.

\begin{figure}[H]
\centering
\resizebox{0.92\textwidth}{!}{%
\begin{ganttchart}[
  hgrid,
  vgrid,
  x unit=0.72cm,
  y unit chart=0.62cm,
  bar height=0.55,
  bar/.append style={fill=azulUC, rounded corners=2pt},
  group/.append style={fill=azulOscuro, draw=azulOscuro},
  milestone/.append style={fill=moradoUC},
  title/.style={fill=azulClaro, draw=azulUC},
  title label font=\tiny\bfseries,
  bar label font=\tiny,
  group label font=\tiny\bfseries,
  milestone label font=\tiny,
  canvas/.append style={draw=azulUC}
]{1}{24}

\gantttitle{Semanas de ejecución}{24} \\
\gantttitlelist{1,2,3,4,5,6,7,8,9,10,11,12,13,14,15,16,17,18,19,20,21,22,23,24}{1} \\

\ganttgroup{Fase 1: Definición}{1}{2} \\
\ganttbar{Diagnóstico del problema}{1}{2} \\
\ganttbar{Formulación de preguntas e hipótesis}{2}{2} \\

\ganttgroup{Fase 2: Revisión de literatura}{2}{5} \\
\ganttbar{Antecedentes y marco teórico}{2}{4} \\
\ganttbar{Normativa y reglas de negocio}{4}{5} \\

\ganttgroup{Fase 3: Diseño metodológico}{6}{8} \\
\ganttbar{Instrumentos y validación}{6}{7} \\
\ganttbar{Definición de muestra}{7}{8} \\

\ganttgroup{Fase 4: Recolección}{8}{12} \\
\ganttbar{Encuestas y entrevistas}{8}{11} \\
\ganttbar{Observación de procesos}{9}{11} \\
\ganttbar{Extracción de datos históricos}{10}{12} \\

\ganttgroup{Fase 5: Análisis}{12}{15} \\
\ganttbar{Codificación cualitativa}{12}{13} \\
\ganttbar{Estadística y preparación del dataset}{13}{15} \\

\ganttgroup{Fase 6: Diseño de la solución}{14}{20} \\
\ganttbar{Requerimientos y modelado UML}{14}{17} \\
\ganttbar{Algoritmo de soluciones rápidas}{17}{20} \\
\ganttbar{Motor analítico y tableros}{18}{20} \\

\ganttgroup{Fase 7: Validación}{19}{22} \\
\ganttbar{Prueba piloto}{19}{21} \\
\ganttbar{Validación con expertos}{21}{22} \\

\ganttgroup{Fase 8: Documentación}{20}{24} \\
\ganttbar{Redacción del informe}{20}{23} \\
\ganttbar{Socialización y cierre}{23}{24} \\

\ganttmilestone{Entrega final}{24}

\end{ganttchart}
}
\caption{Cronograma general del proyecto (24 semanas)}
\label{fig:gantt}
\end{figure}

\section{Plan de Mitigación de Riesgos}

\begin{table}[H]
\centering
\caption{Riesgos, impacto y plan de mitigación}
\label{tab:riesgos}
\small
\begin{tabularx}{\textwidth}{@{}p{1.2cm} X c c X@{}}
\toprule
\rowcolor{naranjaClaro}
\textbf{ID} & \textbf{Riesgo} & \textbf{Prob.} & \textbf{Impacto} & \textbf{Mitigación} \\
\midrule
R-01 & Baja participación en encuestas y evaluación docente. &
  Media & Alto &
  Recordatorios automáticos, tiempo protegido en clase, incentivos
  académicos y reducción de la longitud del instrumento. \\
R-02 & Resistencia al cambio de docentes y personal administrativo. &
  Media & Alto &
  Socialización temprana, capacitación, pilotos con voluntarios y
  visibilización de beneficios. \\
R-03 & Datos históricos incompletos o inconsistentes. &
  Alta & Medio &
  Reglas de limpieza documentadas, imputación controlada y exclusión de
  registros no confiables. \\
R-04 & Percepción de vigilancia sobre la actividad docente. &
  Media & Alto &
  Anonimato estricto, transparencia de criterios y uso exclusivamente
  formativo de los resultados. \\
R-05 & Retraso en la integración con el sistema central. &
  Media & Medio &
  Capa de abstracción (\ac{API} \ac{REST}) y datos simulados para pruebas
  tempranas. \\
R-06 & Bajo desempeño de los modelos analíticos. &
  Media & Medio &
  Validación cruzada, comparación contra línea base y ajuste iterativo
  de variables. \\
R-07 & Incumplimiento de plazos por carga operativa del equipo. &
  Alta & Medio &
  Priorización por valor, entregas incrementales y revisión semanal del
  cronograma. \\
R-08 & Incidente de seguridad o fuga de información. &
  Baja & Muy alto &
  Cifrado, \ac{RBAC}, auditoría de accesos, respaldo y anonimización de
  \ac{PII}. \\
\bottomrule
\end{tabularx}
\end{table}

\section{Recursos Requeridos}

\begin{table}[H]
\centering
\caption{Recursos humanos, técnicos y de infraestructura}
\label{tab:recursos}
\small
\begin{tabularx}{\textwidth}{@{}l X@{}}
\toprule
\rowcolor{verdeClaro}
\textbf{Categoría} & \textbf{Recursos} \\
\midrule
Humanos & Investigador principal, analista de datos, desarrollador \ac{API},
  diseñador de interfaz, asesor metodológico, contraparte de secretaría. \\
Software & Entorno de desarrollo, gestor de base de datos relacional,
  herramienta de analítica, editor de documentación, control de versiones. \\
Infraestructura & Servidor de aplicaciones, servidor de base de datos,
  repositorio analítico, dominio institucional con certificado TLS. \\
Datos & Registros históricos de cupos, materias, faltas, evaluaciones y
  quejas del \ac{DIC}, debidamente anonimizados. \\
\bottomrule
\end{tabularx}
\end{table}

\section{Presupuesto Estimado}

\begin{table}[H]
\centering
\caption{Presupuesto estimado del proyecto}
\label{tab:presupuesto}
\small
\begin{tabularx}{\textwidth}{@{}l X r r@{}}
\toprule
\rowcolor{azulClaro}
\textbf{Rubro} & \textbf{Descripción} & \textbf{Unidades} & \textbf{\%} \\
\midrule
Talento humano & Investigación, análisis de datos y desarrollo & 60 & 60\% \\
Infraestructura & Servidores y almacenamiento (periodo del proyecto) & 15 & 15\% \\
Licencias & Herramientas de analítica y notificaciones & 10 & 10\% \\
Capacitación & Formación de docentes y secretaría & 8 & 8\% \\
Imprevistos & Contingencia del proyecto & 7 & 7\% \\
\midrule
\rowcolor{grisClaro}
\textbf{Total} & & \textbf{100} & \textbf{100\%} \\
\bottomrule
\end{tabularx}
\end{table}

\section{Criterios de Aceptación}

\begin{table}[H]
\centering
\caption{Criterios de aceptación del proyecto}
\label{tab:aceptacion}
\small
\begin{tabularx}{\textwidth}{@{}p{1.4cm} X@{}}
\toprule
\rowcolor{moradoClaro}
\textbf{ID} & \textbf{Criterio verificable} \\
\midrule
CA-01 & Los 56 requerimientos funcionales del Capítulo~\ref{ch:requerimientos}
  están trazados a objetivos estratégicos y casos de uso. \\
CA-02 & Los diagramas \ac{UML} compilan en \LaTeX{} y son consistentes con la
  especificación de requerimientos. \\
CA-03 & El tiempo medio de respuesta a solicitudes de cupo es inferior a 24 horas
  durante la prueba piloto. \\
CA-04 & El 100\% de las evaluaciones docentes del periodo piloto se completan
  dentro de la ventana configurada. \\
CA-05 & El 95\% de las quejas registradas cuentan con responsable, plazo y estado
  de trazabilidad. \\
CA-06 & El modelo de preferencia horaria alcanza una aceptación promedio superior
  a 4.0 sobre 5 en el grupo piloto. \\
CA-07 & El motor analítico genera al menos 10 recomendaciones priorizadas y el
  70\% de ellas son aplicadas por la secretaría o el departamento. \\
CA-08 & El tablero de verificación está disponible para secretaría y
  departamento con todos los paneles definidos. \\
\bottomrule
\end{tabularx}
\end{table}

\section{Conclusiones}

\begin{cajaObjetivo}[Conclusiones del documento]
\begin{enumerate}[leftmargin=1.2cm, itemsep=4pt]
  \item El problema del \ac{DIC} \textbf{no es de capacidad tecnológica sino de
        integración}: los procesos existen, pero operan de manera aislada, sin
        trazabilidad ni analítica, lo que provoca demoras, evaluaciones
        incompletas y decisiones sin evidencia.

  \item Se evidenció que \textbf{once procesos críticos} concentran la mayor
        parte de las dificultades: cupos, materias, comunicación bidireccional,
        faltas, evaluación docente, quejas, efectividad de la secretaría,
        verificación de resultados, asignación de horarios y ausencia de
        analítica.

  \item La metodología \textbf{mixta con diseño de investigación--acción} resulta
        adecuada porque permite comprender las causas humanas del problema y, al
        mismo tiempo, medir objetivamente variables como tiempos de respuesta,
        tasas de participación y frecuencia de quejas.

  \item El uso de \ac{CRISP-DM} ofrece una ruta estructurada y auditable para el
        componente analítico, garantizando que las recomendaciones se evalúen
        contra una línea base antes de desplegarse.

  \item El \textbf{modelado \ac{UML} completo} (casos de uso, clases, secuencia,
        actividades, estados y componentes) constituye una especificación
        suficiente para iniciar la construcción del sistema sin ambigüedades
        funcionales.

  \item El \textbf{modelo de preferencias por probabilidad} convierte los
        cuestionarios de horarios en una herramienta de decisión objetiva,
        reemplazando la programación académica basada en intuición por una
        basada en evidencia.

  \item El \textbf{algoritmo de soluciones rápidas} con priorización por
        impacto, urgencia y esfuerzo permite a la secretaría y al departamento
        actuar sobre lo verdaderamente relevante, y no solo sobre lo más
        visible.

  \item La \textbf{verificación de resultados mediante tableros} cierra el ciclo
        de mejora continua: cada acción queda registrada, medida y disponible
        para la siguiente iteración de decisión.
\end{enumerate}
\end{cajaObjetivo}

\section{Recomendaciones}

\begin{enumerate}[leftmargin=1.4cm, itemsep=4pt]
  \item \textbf{Implementar por incrementos:} iniciar con los módulos de cupos,
        comunicación, faltas y evaluación docente, que concentran el mayor
        impacto sobre la oportunidad.

  \item \textbf{Institucionalizar el porcentaje mínimo de participación:} la
        evaluación docente debe ser requisito para la matrícula del siguiente
        periodo, con anonimato garantizado.

  \item \textbf{Capacitar a la secretaría en el uso analítico:} no basta con
        entregar tableros; se requiere formación en interpretación de
        indicadores y priorización.

  \item \textbf{Publicar los criterios de priorización de cupos:} la
        transparencia es requisito para la legitimidad del sistema ante los
        estudiantes.

  \item \textbf{Revisar periódicamente las reglas del algoritmo:} los umbrales
        (saturación, inasistencia, plazos) deben ajustarse con datos reales de
        cada periodo académico.

  \item \textbf{Diseñar el plan de protección de datos:} clasificar la
        información, definir tiempos de retención y establecer controles de
        acceso antes de la puesta en producción.

  \item \textbf{Extender progresivamente a otras unidades académicas:} la
        arquitectura modular permite replicar el modelo, adaptando reglas de
        negocio y catálogos.

  \item \textbf{Medir el impacto social:} además de los \ac{KPI} operativos, se
        recomienda medir la percepción de equidad, la reducción de conflictos y
        el clima académico percibido.
\end{enumerate}

\section{Trabajo Futuro}

\begin{table}[H]
\centering
\caption{Líneas de trabajo futuro}
\label{tab:trabajo-futuro}
\small
\begin{tabularx}{\textwidth}{@{}p{3.9cm} X c@{}}
\toprule
\rowcolor{azulClaro}
\textbf{Línea} & \textbf{Descripción} & \textbf{Horizonte} \\
\midrule
Asistente conversacional &
  Interfaz de lenguaje natural para consultar cupos, horarios y estados de casos. &
  Corto \\
Recomendación personalizada &
  Sugerencia de materias y horarios según historial y perfil del estudiante. &
  Mediano \\
Detección temprana de deserción &
  Modelos predictivos sobre inasistencia, desempeño y quejas. &
  Mediano \\
Integración total con el sistema central &
  Sincronización bidireccional de matrícula, notas y titulación. &
  Mediano \\
Analítica de aprendizaje &
  Correlación entre desempeño docente y resultados académicos de los estudiantes. &
  Largo \\
Aplicación móvil nativa &
  Notificaciones push y operaciones sin conexión. &
  Largo \\
\bottomrule
\end{tabularx}
\end{table}